\section*{第 \ref{ch:g9:ions} 章——离子与离子溶液}

\begin{solution}{exo:g9:ions:1}
离子是指失去或得到电子并带有电荷的
原子或原子团。阳离子带正电（失去电子），
阴离子带负电（得到电子）。
\end{solution}

\begin{solution}{exo:g9:ions:2}
\ce{Mg^{2+}}：12 个质子，10 个电子。
\end{solution}

\begin{solution}{exo:g9:ions:3}
\ce{KCl}、\ce{MgCl2}、\ce{CaO}。
\end{solution}

\begin{solution}{exo:g9:ions:4}
食盐水中含有离子 \ce{Na+(aq)} 和 \ce{Cl-(aq)}，它们能移动并
传导电荷。蔗糖分子不带电荷。
\end{solution}

\begin{solution}{exo:g9:ions:5}
\ce{CuSO4}、\ce{AlCl3}、\ce{Na2CO3}。
\end{solution}

\begin{solution}{exo:g9:ions:6}
铁（Ⅲ），\ce{Fe^{3+}}：\ce{Fe^{3+}(aq) + 3OH-(aq) -> Fe(OH)3(s)}。
\end{solution}

\begin{solution}{exo:g9:ions:7}
硫酸根：多 2 个电子。硝酸根：多 1 个。
\end{solution}

\begin{solution}{exo:g9:ions:8}
与硝酸银生成白色沉淀的是 \ce{Cl-}。与氢氧化钠
生成蓝色沉淀的是 \ce{Cu^{2+}}。氯化铜（Ⅱ），\ce{CuCl2}。
\end{solution}

\begin{solution}{exo:g9:ions:9}
四个：每侧各一个（左、右、上、下）。
\end{solution}

\begin{solution}{exo:g9:ions:10}
\ce{Fe^{2+}(aq) + 2OH-(aq) -> Fe(OH)2(s)}。
\end{solution}

\begin{solution}{exo:g9:ions:11}
氯化锌：钠离子与氢氧化钠不产生沉淀，
锌离子则生成白色沉淀：
\ce{Zn^{2+}(aq) + 2OH-(aq) -> Zn(OH)2(s)}。
\end{solution}

\begin{solution}{exo:g9:ions:12}
\ce{CaCl2}：2000 个氯离子。电荷总和：$1000 \times (+2) + 2000 \times
(-1) = 0$。
\end{solution}

\begin{solution}{pb:g9:ions:1}
\textbf{1.} 铁（Ⅱ），\ce{Fe^{2+}}：\ce{Fe(OH)2}。

\textbf{2.} 铁（Ⅲ），\ce{Fe^{3+}}：\ce{Fe(OH)3}。

\textbf{3.} \ce{Fe^{2+}} 多 1 个电子：$26 - 2 = 24$，
而 \ce{Fe^{3+}} 为 $26 - 3 = 23$。

\textbf{4.} \ce{Fe^{3+}(aq) + 3OH-(aq) -> Fe(OH)3(s)}。

\textbf{5.} 流出时含有铁（Ⅱ）离子，与空气接触后
转化为铁（Ⅲ）离子，进而形成橙色固体。

\textbf{6.} \ce{Fe(OH)3}：$(+3) + 3 \times (-1) = 0$。

\textbf{7.} 流出时则不然：铁离子已溶解，可穿过
滤纸。橙色固体形成后则可以，它是一种
沉淀，可被滤纸截留。

\textbf{8.} 否：它是混合物（水和溶解的离子）。刚取出时
且澄清，它是均相混合物。

\textbf{9.} $56 \times 1.67 \times 10^{-27} = \qty{9.35e-26}{kg}$。

\textbf{10.} $\qty{0.30}{mg} = \qty{0.30e-6}{kg} = \qty{3.0e-7}{kg}$。

\textbf{11.} 离子与原子相比，相差两个或三个电子，
这些电子的质量与原子核相比可忽略不计。

\textbf{12.} $3.0 \times 10^{-7} / 9.35 \times 10^{-26} \approx 3.2
\times 10^{18}$ 个铁离子每升。
\end{solution}
